\section{案例研究：CartPole 与 PPO}
\subsection{环境与任务介绍}
CartPole 是经典的 balancing 问题：策略需要控制小车使得杆保持直立，reward 通常为 survival time。每条轨迹长度有限，适合 policy gradient 的在线收敛展示。

\subsection{训练流程与日志}
我们用 PPO 进行训练，采样 2048 步，A2C style 可能在本阶段失败。训练日志典型包含：total reward、policy entropy、clip fraction、value loss。
\begin{table}[H]
\centering
\small
\begin{tabular}{p{0.3\textwidth}p{0.5\textwidth}}
\toprule
Metric & Interpretation \\
\midrule
Clip fraction & 探测梯度是否常 hit clip threshold \\
Entropy & policy 变化度 \\
Value loss & Critic 是否收敛 \\
Reward std & 策略的鲁棒性 \\
\bottomrule
\end{tabular}
\caption{CartPole PPO 训练中常用指标}
\end{table}

\subsection{复盘与调优}
如果 reward 停滞，可通过增加 entropy coefficient、减小 learning rate、增加 horizon 进行复盘。每次调优后记录 prompt \& hyper parameters，并附上 `override note` 以便 future engineer 复现。

\begin{knowledgebox}{复盘步骤}
\begin{itemize}
    \item 检查 reward curve 是否 Plateau；
    \item 观察 advantage 分布与 value loss；
\item 调整 entropy / clip threshold，记录 seed；
\item 执行 soft reset（将 critic 权重设置为 previous checkpoint），观察是否恢复。
\end{itemize}
\end{knowledgebox}

\subsection{本章小结}
CartPole + PPO 提供了完整 pipeline：从采样、估计、更新、到日志/复盘。此案例可作为大型 benchmark 的 micro 模型供新人快速掌握 policy gradient training。

